\bookchapter{Chapitre 25 : Traitement des smartphones d'occasion et patrimoine numérique}

\chapterintro{Hé les amis, après plus de vingt chapitres d'« entraînement », vous êtes passé d'un « débutant du numérique » à un utilisateur de smartphone plutôt à l'aise. Mais la vie technologique connaît des cycles — le nouveau devient toujours l'ancien. Quand vous êtes prêt à laisser votre vieux compagnon « prendre sa retraite », savez-vous comment lui dire adieu dignement et en toute sécurité ?}

\booksection{I. « Euthanasie » du vieux smartphone : guide de vente ou de don sécurisé}

Principe fondamental : le nettoyage des données n'est pas une « suppression », mais un « écrasement ». Beaucoup pensent qu'il suffit de glisser les photos dans la corbeille et de vider l'album pour que tout soit réglé, mais après la revente du smartphone, le nouveau propriétaire n'a qu'à lancer un logiciel de récupération de données pour retrouver les anciennes photos, les historiques de discussion et les photos de cartes d'identité — ce type de « fiasco social » se produit chaque année.

\begin{enumerate}
  \item \textbf{Grand ménage personnel}: Effectuez une dernière sauvegarde complète, déconnectez tous les comptes (iCloud, compte Huawei/Xiaomi/Samsung, compte Google), désactivez la fonction « Localiser mon appareil » — sinon le prochain utilisateur sera bloqué par le verrou d'activation.
  \item \textbf{Exécuter le « formatage » ultime}: La « réinitialisation d'usine » des smartphones modernes inclut généralement un effacement chiffré — elle détruit la clé de chiffrement elle-même, une fois la clé détruite, toutes les données du disque se transforment instantanément en données inintelligibles.
  \item \textbf{Dernier « leurre » (uniquement pour les anciens appareils)}: Les smartphones modernes (iOS 8+, Android 6+) sont chiffrés par défaut (FBE), la réinitialisation d'usine détruit la clé de chiffrement, les données sont théorlement irrécupérables,\textbf{cette étape n'est pas nécessaire}. Uniquement pour les très anciens appareils non chiffrés (avant 2015), il est nécessaire de remplir le stockage avec des fichiers poubelle après la réinitialisation d'usine, puis d'effectuer une nouvelle réinitialisation d'usine.
\end{enumerate}

\textbf{Tableau de référence des chemins d'effacement de données par marque}: Bien que le principe soit le même, chaque marque cache légèrement différemment l'emplacement de la « réinitialisation d'usine », les différences clés sont les suivantes.

\begin{center}
\begin{tabularx}{\linewidth}{|X|X|X|}
\hline
\textbf{Marque} & \textbf{Entrée d'effacement} & \textbf{Différences clés et points d'attention} \tabularnewline
\hline
Huawei & Réglages > Système et mises à jour > Réinitialisation > Réinitialisation d'usine & Il faut d'abord se déconnecter du compte Huawei et désactiver « Localiser mon appareil », sinon le verrou d'activation se déclenche ; en mode pur, une vérification par mot de passe peut être nécessaire. \tabularnewline
Xiaomi / Redmi & Réglages > Mon appareil > Réinitialisation d'usine & En bas se trouvent les options « Formater la carte SD virtuelle » et « Effacer tout le contenu et les réglages », les deux doivent être cochées ; avant de se déconnecter du compte Xiaomi, désactivez « Localiser l'appareil ». \tabularnewline
OPPO / OnePlus & Réglages > Autres réglages > Restaurer le smartphone > Effacer complètement toutes les données & ColorOS propose un mode « effacement complet » qui couvre plusieurs fois le stockage ; la saisie du code de déverrouillage est requise pour confirmer. \tabularnewline
vivo & Réglages > Gestion du système > Sauvegarde et réinitialisation > Effacer toutes les données & Cochez « Formater le stockage du smartphone » pour nettoyer la zone de la carte SD ; il faut se déconnecter du compte vivo et désactiver « Localiser le smartphone ». \tabularnewline
Samsung & Réglages > Gestion générale > Réinitialisation > Réinitialisation d'usine & Il faut d'abord supprimer le compte Samsung et désactiver « Localiser mon appareil » ; certains modèles ont une option « écrasement automatique », dont l'activation est recommandée. \tabularnewline
iPhone & Réglages > Général > Transférer ou réinitialiser l'iPhone > Effacer tout le contenu et les réglages & Il faut impérativement se déconnecter de l'Apple ID au préalable, sinon le prochain utilisateur sera bloqué par le verrou d'activation iCloud ; iOS 15+ prend en charge « conserver les données iCloud lors de la déconnexion » pour consultation ultérieure. \tabularnewline
\hline
\end{tabularx}
\end{center}

\begin{tcbwarn}{Les trois choses les plus facilement oubliées}
\begin{itemize}
  \item \textbf{Compte non déconnecté}: Le smartphone est verrouillé à distance après la vente, et vous ne pouvez plus y accéder non plus, vous devez finalement vous rendre au service après-vente officiel avec votre pièce d'identité pour le faire déverrouiller.
  \item \textbf{Carte SD non retirée}: La réinitialisation d'usine ne formate pas les photos de la carte SD externe, n'oubliez pas de la retirer manuellement ou de la formater séparément.
  \item \textbf{Clé USB / périphérique OTG non débranché}: Si un périphérique est connecté pendant l'effacement, les données du périphérique peuvent être effacées par erreur, ou l'effacement peut être interrompu.
\end{itemize}

\end{tcbwarn}

\booksection{II. Patrimoine numérique : comment placer votre âme sociale ?}

Les données dans le smartphone ne sont plus « virtuelles » depuis longtemps. Elles portent votre argent, vos relations sociales, vos créations et même votre « identité numérique ».\textbf{Après le décès d'une personne, comment gérer ces actifs ?}Ce n'est plus une question abstraite, mais un dilemme réel auquel de plus en plus de familles sont confrontées — les proches ne connaissent pas les mots de passe, les plateformes ne permettent pas l'héritage, les cryptomonnaies sont définitivement perdues.

\begin{tcbinfo}{Base juridique : article 49 de la loi sur la protection des informations personnelles}
L'article 49 de la « Loi de la République populaire de Chine sur la protection des informations personnelles », entrée en vigueur en 2021, stipule clairement :\textbf{En cas de décès d'une personne physique, ses proches parents peuvent, pour leurs propres intérêts légitimes et légaux, exercer les droits de consultation, de copie, de correction et de suppression des informations personnelles du défunt tels que définis dans ce chapitre ; sauf arrangement contraire pris par le défunt de son vivant.}

Cela signifie que :

\begin{itemize}
  \item \textbf{Les proches parents ont une base juridique}pour demander aux plateformes de traiter les comptes et les informations personnelles du défunt.
  \item Les « intérêts légitimes et légaux » incluent généralement la gestion de la succession, la fermeture de comptes et la suppression d'informations privées sensibles.
  \item Si le défunt a pris des dispositions par testament ou accord de son vivant (par exemple, désigner un contact de succession), ses dispositions priment.
  \item Les plateformes ne peuvent pas refuser catégoriquement de coopérer sous prétexte de « vie privée de l'utilisateur », mais les procédures spécifiques varient selon les plateformes.
\end{itemize}

\end{tcbinfo}

\booksubsection{WeChat et Alipay : le compte ne peut pas être hérité, mais les fonds le peuvent}

Ces deux applications les plus utilisées par les Chinois ont des règles différentes :

\begin{itemize}
  \item \textbf{Compte WeChat}: Selon le « Contrat de licence et de service du logiciel Tencent WeChat », la propriété du compte appartient à Tencent, l'utilisateur ne dispose que d'un droit d'usage,\textbf{le compte lui-même ne peut pas être hérité ni transféré}. Mais les\textbf{fonds du solde WeChat et du Qianlibao}appartiennent au patrimoine personnel, les proches parents peuvent demander à Tenpay le retrait ou le transfert sur présentation d'un certificat de décès, d'un certificat de parenté et d'un acte notarié de succession.
  \item \textbf{Alipay}: Le compte ne peut pas non plus être hérité directement, mais le solde Alipay, Yu'ebao, Licaitong, etc.\textbf{les fonds sont un patrimoine successible}. Les héritiers peuvent procéder via la « ligne d'assistance client » d'Alipay ou en soumettant hors ligne un certificat de décès et un acte notarié de succession ; les produits de gestion de patrimoine d'Ant Fortune doivent passer par la procédure de la société de fonds.
  \item \textbf{Historique des discussions}: Juridiquement, cela relève de la vie privée personnelle, la plateforme ne le fournira pas spontanément à la famille. S'il y a des indices de succession, ils peuvent être obtenus par voie judiciaire ou notariée.
\end{itemize}

\begin{tcbwarn}{Ce qu'il ne faut absolument pas faire}
\begin{itemize}
  \item \textbf{Ne remettez pas directement les mots de passe du compte du défunt à la famille pour « continuer à l'utiliser »}— cela enfreint le contrat de la plateforme, et en cas de gel des fonds ou de déclenchement du contrôle des risques, la charge de la preuve est encore plus difficile.
  \item \textbf{N'essayez pas d'utiliser le smartphone du défunt pour recevoir des codes de vérification et « retirer des fonds »}— cela peut être qualifié d'usurpation et constituer une infraction pénale.
  \item \textbf{Ne traînez pas}: Certains comptes sont automatiquement annulés après une longue période d'inactivité, ce qui complique la récupération des fonds.
\end{itemize}

\end{tcbwarn}

\booksubsection{« Contact de succession » d'Apple : désignez-le à l'avance, un effort minime pour un grand bénéfice}

Depuis iOS 15.2, Apple a lancé la fonction « Contact de succession », permettant aux utilisateurs de désigner à l'avance un contact de confiance ; après le décès, ce contact peut, sur présentation d'\textbf{un certificat de décès + la clé d'accès du contact de succession}demander l'accès aux données iCloud du défunt (photos, notes, contacts, mémos, etc.). Étapes :

\begin{enumerate}
  \item Sur l'iPhone, accédez à\textbf{[Réglages] > [Votre nom] > [Connexion et sécurité] > [Contact de succession]}.
  \item Touchez « Ajouter un contact de succession », confirmez votre identité avec Face ID.
  \item Choisissez un contact (sélectionnez dans les contacts), le système générera une\textbf{clé d'accès}.
  \item Transmettez la clé d'accès par SMS, e-mail ou en l'imprimant sur papier, et rappelez à l'autre personne de la conserver soigneusement.
  \item Vous pouvez ajouter plusieurs contacts comme sauvegarde, pour éviter un point de défaillance unique.
\end{enumerate}

La personne désignée n'a pas besoin d'un appareil Apple, elle peut demander l'accès sur le site d'Apple avec la clé et le certificat de décès. Le processus est généralement terminé en 7 jours.

\booksubsection{« Gestionnaire de compte inactif » de Google : gère automatiquement vos affaires posthumes}

Google propose le « Gestionnaire de compte inactif (Inactive Account Manager) », qui peut exécuter automatiquement des actions prédéfinies en cas d'inactivité prolongée de l'utilisateur — c'est actuellement l'un des mécanismes de patrimoine numérique les plus aboutis.

\begin{enumerate}
  \item Visitez\textbf{myaccount.google.com/inactive}, connectez-vous à votre compte Google.
  \item Définissez le délai d'inactivité : déclenchement après 3 / 6 / 12 / 18 mois sans connexion.
  \item Ajoutez des contacts de confiance (10 maximum), et définissez les données accessibles à chacun (Gmail, Photos, Drive, etc.).
  \item Optionnel : configurez un e-mail de réponse automatique informant les correspondants que vous ne pouvez plus traiter les messages.
  \item Dernier choix : à l'expiration, faut-il\textbf{partager les données avec les contacts}ou\textbf{supprimer directement le compte}.
  \item Définissez un numéro de téléphone comme vérification de secours pour éviter un déclenchement par erreur.
\end{enumerate}

\booksubsection{Comptes de réseaux sociaux nationaux : commémoration et procédure de suppression}

Les principales plateformes de réseaux sociaux nationales traitent les comptes des défunts selon deux approches :\textbf{« Commémoration » (conserver le compte pour que les proches puissent se recueillir)}et\textbf{« Suppression » (effacement définitif)}. La famille peut demander au service client de la plateforme sur présentation d'un certificat de décès, d'un certificat de parenté et de la pièce d'identité du demandeur.

\begin{itemize}
  \item \textbf{Weibo}: Prend en charge la « commémoration » du compte, sur demande de la famille avec certificat de décès, le compte est verrouillé et ne peut plus être connecté mais le contenu est conservé, la page d'accueil affiche « Ce compte a été défini comme compte commémoratif ». La suppression peut également être demandée, un acte notarié de succession est requis.
  \item \textbf{Douyin}: Via « Moi - Réglages - Aide et commentaires - Compte et sécurité - Suppression de compte » ou le service client, soumettez le certificat de décès, le certificat de parenté et la pièce d'identité du demandeur ; après approbation, le compte peut être supprimé ou converti en statut commémoratif.
  \item \textbf{Bilibili}: Soumettez via le canal du service client le certificat du défunt et le certificat de parenté, vous pouvez demander la suppression du compte ou sa conservation en tant que page commémorative ; le traitement des vidéos historiques des créateurs doit tenir compte des déclarations de droits d'auteur.
  \item \textbf{Xiaohongshu}: Via le service client officiel, soumettez le certificat de décès et le certificat de parenté, vous pouvez demander la suppression du compte et l'effacement du contenu.
\end{itemize}

\begin{tcbinfo}{Rappel important}
\begin{itemize}
  \item Les règles de chaque plateforme peuvent évoluer, consultez les dernières instructions du service client officiel avant de procéder.
  \item Il est recommandé de numériser et de sauvegarder tous les documents justificatifs avant de les soumettre, et de conserver soigneusement les originaux.
  \item Pour un traitement multiplateforme en une seule fois, vous pouvez demander un seul\textbf{acte notarié de succession}, pour éviter des démarches répétées.
\end{itemize}

\end{tcbinfo}

\booksubsection{Le défi de l'héritage des actifs numériques}

Outre les comptes de réseaux sociaux, de plus en plus de personnes détiennent des\textbf{actifs numériques réellement précieux}, dont l'héritage est plus complexe :

\begin{itemize}
  \item \textbf{Cryptomonnaies (Bitcoin, Ethereum, etc.)}: Basé sur le principe « la clé privée est la propriété »,\textbf{sans la clé privée, les actifs disparaissent}. Les héritiers ne peuvent pas « faire appel » sur la chaîne via des documents juridiques. Recommandation : stockez la phrase de récupération / la clé privée fragmentée dans plusieurs emplacements physiques sécurisés (par exemple coffre-fort bancaire + coffre-fort à domicile) et indiquez le mode d'accès dans le testament.
  \item \textbf{Comptes d'échange}: Similaires aux comptes bancaires, vous pouvez demander la disposition auprès de l'échange (par exemple Binance, OKX) avec un acte notarié de succession, mais le processus d'un échange transfrontalier est long.
  \item \textbf{Noms de domaine}: Patrimoine successible, les héritiers peuvent effectuer le transfert auprès du registraire de noms de domaine (par exemple Alibaba Cloud, GoDaddy) après soumission d'un acte notarié de succession.
  \item \textbf{Propriété intellectuelle (articles originaux, vidéos, brevets, droits d'auteur logiciels)}: Les droits patrimoniaux du droit d'auteur sont successibles (les droits moraux ne le sont pas), le droit aux revenus peut être exercé par les héritiers jusqu'à 50 ans après le décès de l'auteur.
  \item \textbf{Stockage cloud et disques réseau}: Les données de contenu appartiennent à la personne et sont successibles ; mais certains contrats de plateforme stipulent que le compte lui-même n'est pas transférable, il faut d'abord télécharger puis supprimer.
\end{itemize}

\textbf{Le coût réel de la « disparition » des cryptomonnaies}\par

Selon un rapport de Chainalysis, plus de\textbf{3 millions de bitcoins}sont déjà perdus à jamais dans le monde en raison de clés privées perdues, soit un sixième du total existant. Si vous détenez des actifs cryptographiques, assurez-vous qu'au moins une personne en qui vous avez une confiance absolue connaisse le\textbf{chemin de récupération}— sinon, ces richesses resteront endormies pour toujours sur la blockchain.

\booksubsection{Modèle de « testament numérique »}

Le testament patrimonial traditionnel ne peut pas gérer les actifs numériques. Il est recommandé de mettre à jour chaque année une annexe de « testament numérique », à conserver avec le testament officiel. Le modèle suivant est fourni à titre indicatif :

\textbf{« Annexe de disposition des actifs numériques personnels » (modèle)}\par

\textbf{Partie 1 : Liste des comptes}(Enregistrez uniquement le nom du compte et la plateforme,\textbf{ne notez pas les mots de passe})

\begin{itemize}
  \item Réseaux sociaux : WeChat (compte xxx), Weibo, Douyin, Bilibili, Xiaohongshu, etc.
  \item Paiement : Alipay, WeChat Pay, applications de cartes bancaires
  \item Cloud : iCloud, Huawei Cloud, Xiaomi Cloud, compte Google, compte Microsoft
  \item E-mail : QQ Mail, NetEase Mail, Gmail, etc.
  \item Autres : disques réseau, services d'abonnement, comptes de jeux, etc.
\end{itemize}

\textbf{Partie 2 : Solution de gestion des mots de passe}

\begin{itemize}
  \item Méthode de stockage du mot de passe principal : utilisez un gestionnaire de mots de passe (par exemple Bitwarden, 1Password), écrivez le mot de passe principal sur papier et conservez-le scellé dans un coffre-fort à domicile.
  \item Clé de récupération / phrase de récupération : stockage fragmenté, au moins deux emplacements physiques (par exemple coffre bancaire + coffre-fort à domicile).
  \item Personne autorisée à consulter : nom, numéro de pièce d'identité, relation avec la personne, coordonnées.
\end{itemize}

\textbf{Partie 3 : Volonté de disposition}

\begin{itemize}
  \item Pour chaque type de compte, souhaitez-vous la « commémoration » ou la « suppression ».
  \item Les données iCloud / Google doivent-elles désigner un « contact de succession » ou activer le « gestionnaire de compte inactif ».
  \item Qui hérite des actifs cryptographiques, description du chemin d'accès à la clé privée.
  \item Qui bénéficie des revenus des œuvres soumises au droit d'auteur.
\end{itemize}

\textbf{Partie 4 : Contacts et plateformes}

\begin{itemize}
  \item Nom de l'exécuteur de confiance, coordonnées, contact de secours.
  \item Entrée du service client des plateformes habituelles et liste des documents justificatifs requis.
  \item Date de mise à jour de cette annexe : \_\_\_\_/\_\_/\_\_, date de la prochaine révision : \_\_\_\_/\_\_/\_\_.
\end{itemize}

\begin{tcbwarn}{Avertissement juridique concernant les « mots de passe »}
Écrire les mots de passe directement dans le testament présente un risque de fuite, et après la publication du testament, tous les héritiers pourront les voir.\textbf{Une approche plus sûre consiste à : n'inscrire dans le testament que la « liste des comptes + volonté de disposition + emplacement de gestion des mots de passe », et à transmettre séparément le mot de passe principal du gestionnaire de mots de passe à l'exécuteur de confiance}. Cela répond aux exigences juridiques tout en réduisant le risque de fuite de mots de passe.

\end{tcbwarn}

\booksection{III. Le dernier voyage écologique}

Un smartphone contient des dizaines de métaux tels que l'or, l'argent, le cuivre et les terres rares. Le jeter n'importe où pollue les sols et gaspille des ressources. Il existe trois bons canaux de recyclage, chacun avec ses avantages :

\begin{center}
\begin{tabularx}{\linewidth}{|X|X|X|X|}
\hline
\textbf{Canal} & \textbf{Prix de reprise} & \textbf{Sécurité des données} & \textbf{Scénario d'utilisation} \tabularnewline
\hline
Reprise officielle (Huawei/Apple/Xiaomi, etc.) & Moyen, souvent accompagné de bons de réduction & Le plus élevé, processus officiel normalisé & Le plus avantageux quand vous envisagez de changer pour la même marque \tabularnewline
Plateformes de recyclage tierces (Aihuishou, Zhuanzhuan, Huishoubao, etc.) & Légèrement supérieur, possibilité de comparer les prix & Assez élevé, avec rapport d'effacement des données & Si vous souhaitez encaisser directement sans dépendre d'une marque \tabularnewline
Don caritatif (écoles rurales, associations, etc.) & Aucun retour financier & Effacement complet à faire soi-même & L'appareil est encore utilisable et vous voulez faire une bonne action \tabularnewline
Réutilisation (transformation personnelle) & Économise le coût d'un nouvel appareil & Le propriétaire d'origine continue à l'utiliser, sans risque & Les performances de l'appareil sont correctes, et vous avez un besoin de transformation à la maison \tabularnewline
\hline
\end{tabularx}
\end{center}

\begin{itemize}
  \item \textbf{Reprise/échange officiel}: La méthode la plus recommandée, processus normalisé et sécurité des données garantie. Les sites officiels d'Apple, Huawei, Xiaomi, etc. proposent une entrée « reprise », l'estimation peut déduire du prix d'un nouvel appareil.
  \item \textbf{Plateforme de reprise agréée}: Choisissez une grande plateforme de recyclage réputée ; après commande, SF Express vient récupérer l'appareil à domicile, et le paiement est effectué après inspection. Assurez-vous de choisir\textbf{une plateforme qui fournit un « rapport d'effacement des données »}et conservez le justificatif.
  \item \textbf{Don caritatif}: Vous pouvez contacter des écoles rurales locales ou des organisations caritatives communautaires ; les anciens smartphones peuvent être transformés en tablettes éducatives ou en téléphones pour seniors.\textbf{Avant le don, effectuez impérativement un effacement complet et retirez toutes les cartes SIM/SD}.
  \item \textbf{Seconde vie}: Transformez-le en navigateur de voiture, en centrale domotique, en caméra de surveillance, en babyphone ou en cadre photo numérique. Un ancien smartphone peut encore servir pendant des années.
\end{itemize}

\begin{tcbinfo}{La vérité sur les déchets électroniques}
La carte de circuit imprimé d'un smartphone contient\textbf{de l'or, de l'argent, du palladium et du cuivre}et autres métaux précieux ; chaque tonne de cartes de circuit imprimé usagées contient des dizaines de fois plus d'or que le minerai aurifère. Mais en cas de mise en décharge sauvage, le plomb, le mercure et le cadmium qu'elles contiennent s'infiltrent dans les sols et les nappes phréatiques, polluant pendant des décennies.\textbf{Un recyclage correct protège l'environnement et permet la régénération des métaux rares.}

\end{tcbinfo}

\begin{tcbpractice}{Pratique : organisez des « funérailles standard » pour votre vieux smartphone}
\begin{enumerate}
  \item Dernière sauvegarde complète : photos téléchargées dans le cloud, fichiers importants copiés sur l'ordinateur, historique des discussions exporté et archivé.
  \item Déconnexion de tous les comptes : iCloud / Huawei / Xiaomi / Google, et désactivation de « Localiser mon appareil ».
  \item Retirez les cartes SIM et SD, débranchez tous les périphériques OTG.
  \item Effectuez une réinitialisation d'usine chiffrée, en cochant manuellement « Formater le stockage » si nécessaire.
  \item (Anciens appareils uniquement) Remplissez le stockage avec des fichiers poubelle, puis effectuez à nouveau une réinitialisation d'usine.
  \item Choisissez sa prochaine destination : reprise officielle, plateforme de recyclage agréée, don caritatif, ou don à un membre de la famille.
  \item Mettez à jour votre « inventaire du patrimoine numérique » en rayant les appareils mis au rebut.
\end{enumerate}

\end{tcbpractice}

\begin{tcbpractice}{Liste des tâches du mois}
\begin{enumerate}
  \item Désignez un « contact de succession » Apple ou configurez le « gestionnaire de compte inactif » Google pour votre smartphone principal actuel.
  \item Rédigez une « Annexe de disposition des actifs numériques personnels » à l'aide du modèle de cette section, et communiquez l'emplacement de stockage à votre famille.
  \item Vérifiez s'il y a 2 smartphones ou plus inutilisés à la maison, et planifiez une collecte par le fabricant ou une plateforme agréée.
\end{enumerate}

\end{tcbpractice}
